% Resumo: ambiente abstract e palavras-chave, como no template da disciplina (src/abstract.tex no Overleaf).
% Resultado só qualitativo, sem números (pedido do autor); termina no último resultado, sem frase de implicação.

\begin{abstract}
Modelos de linguagem permitem fazer perguntas em linguagem natural sobre dados de eventos de futebol, mas os dados
da partida precisam chegar ao contexto do modelo. Este trabalho transforma os eventos da final da Copa do Mundo de
2022 em uma tabela de ações e em um grafo, no qual cada ação liga um jogador à partida, e compara \nArms{}
abordagens de aprendizado em contexto para dar esses dados ao modelo: sem dados, a tabela de ações inteira no
prompt, estatísticas resumidas, busca vetorial, consulta escrita pelo modelo na linguagem do grafo (Cypher) e
ferramentas com consultas parametrizadas sobre o grafo. A avaliação usou \nQuestions{} perguntas sobre a partida,
em \nModels{} modelos, dois menores e dois maiores. As ferramentas sobre o grafo foram a única abordagem com acerto
alto em todos os modelos. A consulta escrita pelo modelo só alcançou esse nível nos modelos maiores, com raciocínio estendido. Nesses modelos, a tabela inteira no prompt também chegou perto, mas gastou mais tokens de saída e mais tempo por pergunta. A busca vetorial e as estatísticas resumidas acertaram pouco em todos os modelos.
\end{abstract}

\noindent\textbf{Palavras-chave:} modelos de linguagem, aprendizado em contexto, recuperação de dados, grafos,
dados de eventos de futebol.
